\documentclass[10pt,a4paper]{amsart}
\usepackage{amsmath,amssymb,amsthm,graphicx}
\usepackage{verbatim}
\newenvironment{reference}{\comment}{\endcomment}
%\newenvironment{reference}{}{}
\newenvironment{slogan}{\comment}{\endcomment}
\newenvironment{history}{\comment}{\endcomment}

% For commutative diagrams we use Xy-pic
\usepackage[all]{xy}

% We use 2cell for 2-commutative diagrams.
\xyoption{2cell}
% Theorem environments.
%
\theoremstyle{plain}
\newtheorem{theorem}[subsection]{Theorem}
\newtheorem{proposition}[subsection]{Proposition}
\newtheorem{lemma}[subsection]{Lemma}

\theoremstyle{definition}
\newtheorem{definition}[subsection]{Definition}
\newtheorem{example}[subsection]{Example}
\newtheorem{exercise}[subsection]{Exercise}
\newtheorem{situation}[subsection]{Situation}

\theoremstyle{remark}
\newtheorem{remark}[subsection]{Remark}
\newtheorem{remarks}[subsection]{Remarks}

\numberwithin{equation}{subsection}

% Macros
%
\def\lim{\mathop{\mathrm{lim}}\nolimits}
\def\colim{\mathop{\mathrm{colim}}\nolimits}
\def\Spec{\mathop{\mathrm{Spec}}}
\def\Hom{\mathop{\mathrm{Hom}}\nolimits}
\def\Ext{\mathop{\mathrm{Ext}}\nolimits}
\def\SheafHom{\mathop{\mathcal{H}\!\mathit{om}}\nolimits}
\def\SheafExt{\mathop{\mathcal{E}\!\mathit{xt}}\nolimits}
\def\Sch{\mathit{Sch}}
\def\Mor{\mathop{\mathrm{Mor}}\nolimits}
\def\Ob{\mathop{\mathrm{Ob}}\nolimits}
\def\Sh{\mathop{\mathit{Sh}}\nolimits}
\def\NL{\mathop{N\!L}\nolimits}
\def\CH{\mathop{\mathrm{CH}}\nolimits}
\def\proetale{{pro\text{-}\acute{e}tale}}
\def\etale{{\acute{e}tale}}
\def\QCoh{\mathit{QCoh}}
\def\Ker{\mathop{\mathrm{Ker}}}
\def\Im{\mathop{\mathrm{Im}}}
\def\Coker{\mathop{\mathrm{Coker}}}
\def\Coim{\mathop{\mathrm{Coim}}}

% Boxtimes
%
\DeclareMathSymbol{\boxtimes}{\mathbin}{AMSa}{"02}

%
% Macros for moduli stacks/spaces
%
\def\QCohstack{\mathcal{QC}\!\mathit{oh}}
\def\Cohstack{\mathcal{C}\!\mathit{oh}}
\def\Spacesstack{\mathcal{S}\!\mathit{paces}}
\def\Quotfunctor{\mathrm{Quot}}
\def\Hilbfunctor{\mathrm{Hilb}}
\def\Curvesstack{\mathcal{C}\!\mathit{urves}}
\def\Polarizedstack{\mathcal{P}\!\mathit{olarized}}
\def\Complexesstack{\mathcal{C}\!\mathit{omplexes}}
% \Pic is the operator that assigns to X its picard group, usage \Pic(X)
% \Picardstack_{X/B} denotes the Picard stack of X over B
% \Picardfunctor_{X/B} denotes the Picard functor of X over B
\def\Pic{\mathop{\mathrm{Pic}}\nolimits}
\def\Picardstack{\mathcal{P}\!\mathit{ic}}
\def\Picardfunctor{\mathrm{Pic}}
\def\Deformationcategory{\mathcal{D}\!\mathit{ef}}
\usepackage{cleveref}
\crefname{originallemma}{original lemma}{original lemmas}
\Crefname{originallemma}{Original lemma}{Original lemmas}
\makeatletter
\newlabel{algebra-lemma-injective-minimal-primes-in-image}{{00FK}{source}{}{sourcepointer-1}{}}
\newlabel{algebra-lemma-injective-minimal-primes-in-image@cref}{{[originallemma][0][]00FK}{[1][0][]source}{}{}{}}
\makeatother
\crefname{originallemma}{original lemma}{original lemmas}
\Crefname{originallemma}{Original lemma}{Original lemmas}
\makeatletter
\newlabel{lemma-properties-normalization}{{0ABP}{source}{}{sourcepointer-2}{}}
\newlabel{lemma-properties-normalization@cref}{{[originallemma][0][]0ABP}{[1][0][]source}{}{}{}}
\makeatother
\crefname{originaldefinition}{original definition}{original definitions}
\Crefname{originaldefinition}{Original definition}{Original definitions}
\makeatletter
\newlabel{spaces-properties-definition-dimension-local-ring}{{04NA}{source}{}{sourcepointer-3}{}}
\newlabel{spaces-properties-definition-dimension-local-ring@cref}{{[originaldefinition][0][]04NA}{[1][0][]source}{}{}{}}
\makeatother
\title{Codimension zero points of a normalization are images of codimension zero points}
\author{The Stacks Project Authors}
\date{stacks-project:spaces-morphisms.tex:local-snapshot-20261012; GNU FDL 1.2 or later}
\begin{document}
\maketitle
\setcounter{subsection}{1}
\noindent\textit{Editorial note.} All delivered mathematical text is the author's own, copied from the pinned Stacks Project chapter file spaces-morphisms.tex; no statement or proof step was rewritten, reordered or removed.\par\medskip
\noindent\textit{Editorial note.} The proof uses the codimension zero definition through the local ring dimension of the chapter; the identification of minimal primes with codimension zero points is made explicitly.\par\medskip
\noindent\textit{Editorial note.} The three references of the counted spans are recorded as non-load-bearing navigation pointers; the argument that consumes them stays inside the counted proof span.\par\medskip
\noindent\textit{Editorial note.} Preamble: the three native excerpt spans of sources/stacks-spaces-morphisms/preamble.tex (lines 8-13, 14-18 and 154-217). The chapter file itself carries no preamble text.\par\medskip
\noindent\textit{Allowed standard prerequisites (editorial).}\begin{itemize}
\item Etale locality of the normalization and its affine description. Lemma lemma-properties-normalization, used at proof lines 10118-10120. Reduces to X = Spec(A) and identifies X' = Spec(A') with A' the integral closure in the global sections.
\item Codimension zero points of an affine scheme correspond to minimal primes. Properties of Spaces, Definition spaces-properties-definition-dimension-local-ring, used at proof lines 10111-10120. Turns the point x' into a minimal prime of A'.
\item Existence of an affine etale cover of Y. used at proof lines 10121-10122. Produces V = Spec(B) with B the source of an injective map from A'.
\item The image of a minimal prime along an injective ring map. Algebra, Lemma algebra-lemma-injective-minimal-primes-in-image, used at proof lines 10123-10125. Shows that every minimal prime of A' comes from a minimal prime of B.
\end{itemize}\medskip
\ifdefined\hypertarget\hypertarget{sourcepointer-1}{}\fi
\noindent\textit{Source reference.} Original lemma 00FK: focused-seeds/AGP/sources/stacks-snapshot-20261011/tags.txt, source lines 574--574. Navigation pointer into the Stacks Project; the standard fact it supplies is applied inside the counted proof span and the argument resting on it stays there.\par\medskip
\ifdefined\hypertarget\hypertarget{sourcepointer-2}{}\fi
\noindent\textit{Source reference.} Original lemma 0ABP: focused-seeds/AGP/sources/stacks-snapshot-20261011/tags.txt, source lines 13027--13027. Navigation pointer into the Stacks Project; the standard fact it supplies is applied inside the counted proof span and the argument resting on it stays there.\par\medskip
\ifdefined\hypertarget\hypertarget{sourcepointer-3}{}\fi
\noindent\textit{Source reference.} Original definition 04NA: focused-seeds/AGP/sources/stacks-snapshot-20261011/tags.txt, source lines 6071--6071. Navigation pointer into the Stacks Project; the standard fact it supplies is applied inside the counted proof span and the argument resting on it stays there.\par\medskip
\begin{lemma}
\label{lemma-normalization-generic}
Let $S$ be a scheme. Let $f : Y \to X$ be a quasi-compact and quasi-separated
morphism of schemes. Let $X' \to X$ be the normalization of $X$ in $Y$.
If $x' \in |X'|$ is a point of codimension $0$
(Properties of Spaces, Definition
\ref{spaces-properties-definition-dimension-local-ring}), then
$x'$ is the image of some $y \in |Y|$ of codimension $0$.
\end{lemma}

\begin{proof}
By Lemma \ref{lemma-properties-normalization} and the definitions, we may
assume that $X = \Spec(A)$ is affine. Then $X' = \Spec(A')$ where $A'$ is
the integral closure of $A$ in $\Gamma(Y, \mathcal{O}_Y)$ and $x'$ corresponds
to a minimal prime of $A'$. Choose a surjective \'etale
morphism $V \to Y$ where $V = \Spec(B)$ is affine. Then
$A' \to B$ is injective, hence every minimal prime of $A'$ is
the image of a minimal prime of $B$, see
Algebra, Lemma \ref{algebra-lemma-injective-minimal-primes-in-image}.
The lemma follows.
\end{proof}

\end{document}
